\scanpage{17}
stable par la stratification envisagée. De $=$, comme $VF=0$, $V$ est nul sur $H'$,\inserted{(étant $=\operatorname{Ker}V=\operatorname{Im}F=H'$)} et son image donc contenue dans $H^{10(p)}=\operatorname{Ker}F$, \deleted{car $FV=0$}, et $=$ égale; on trouve que \uline{$H^{10(p)}$ est stable par \deleted{cette}\inserted{la} stratification}. \deleted{donc}

\uline{Attention}, cela \uline{n'implique pas à priori que $H^{10}$ le soit également}. Tout ce qu'on peut dire (valable aussi pour $\Lambda$ un \unsure{anneau} local noeth. compl., corps résiduel $k\ldots$) c'est que si $S=\operatorname{Specf}(\Lambda)$, et si on considère les deux images inverses $E_1,E_2$ de $H$ sur \deleted{$\unsure{S}$} $T=\operatorname{Specf}(\Lambda\widehat\otimes\Lambda)=S\times S$, \deleted{alors} et les deux images inverses $F_1,F_2$ de $H^{10}$, et considérant l'isom de stratification
\diagram{\begin{tikzcd}[column sep=large]
u:E_1\arrow[r,"\sim"] & E_2
\end{tikzcd}}
d'où
\diagram{\begin{tikzcd}[column sep=large]
u^{(p)}:E_1^{(p)}\arrow[r,dash] & E_2^{(p)},
\end{tikzcd}}
alors $u^{(p)}(F_1^{(p)})=\deleted{\unsure{u^{(p)}(}}F_2^{(p)}$, i.e. $u(F_1)^{(p)}=F_2^{(p)}$. Si donc $u(F_1)=G$, $F_2=G'$, $E=E_2$, on aura $G^{(p)}=G'^{(p)}$ comme \deleted{\unsure{éléments}}\inserted{sections de} de $\operatorname{Grass}(E)$ sur $T$, ou ce qui revient au $=$, ces sections deviennent égales par changt de base
\diagram{\begin{tikzcd}[column sep=large,row sep=small]
T'\arrow[r,"\pi","\text{puissance p.ième}"']\arrow[d,equal] & T. \\
T &
\end{tikzcd}}
On ne peut en général en conclure $G=G'$, car en général $\pi$ n'est pas un épim. Mais si \inserted{$S$ donc} $T$ est formellement lisse sur $k$,
\mspage{16}
